\chapter{User Manual and Quick Start Guide}
\section{Installation and Release Selection}
Install Git and Python 3.11 or later. These Windows commands use \file{py};
use your Python executable elsewhere. No third-party runtime packages are needed.
\begin{lstlisting}
git clone https://github.com/Lcurtis1Broncos/CSC-481-Capstone-Group.git
cd CSC-481-Capstone-Group
git checkout v1.0-final
py -m unittest discover -s tests -v
py tests/verify_submission.py submission_check
\end{lstlisting}
The acceptance script generates inputs, invokes the CLIs, and saves outputs
and reports. All seven checks should pass. Always choose a new output folder.

\section{Inspecting the Generated Image}
To create only the portable dataset and inspect its records:
\begin{lstlisting}
py tests/generate_submission_dataset.py generated_dataset
py code/ntfs_image_reader.py generated_dataset/synthetic.dd --list-files 8 --show-record 1 --show-record 2
py code/ntfs_recovery.py generated_dataset/synthetic.dd --record 2 --stream CapstoneEvidence --output new_ads.bin
\end{lstlisting}
The final command exports and hashes the ADS. The acceptance script adds
original expectations and a saved report to establish a verified PASS.

\section{Reproducing the Report Campaign}
From the report branch, which adds the campaign after the code tag, run:
\begin{lstlisting}
py tests/run_report_campaign.py report_campaign_new
\end{lstlisting}
Expected results are 50 exact recoveries, 50 baseline matches, and 50
single-bit hash failures. Originals, paired inputs, reports, and CSV results
are saved in the new folder.

For the separate detection experiment on the report branch:
\begin{lstlisting}
py tests/run_detection_campaign.py detection_campaign_new
\end{lstlisting}
This creates 100 labeled inputs, a manifest, \file{results.csv}, and
\file{summary.json}. To repeat the two Windows positive checks as well,
add \file{--lab-image} followed by the preserved working-image path.
This optional step hashes the entire large image twice. Recorded results,
the 73-test log, a result figure, and reproduction details are retained under
\file{reports/evidence/final_detection/}. The earlier code tag does not
contain either report campaign.

\section{Historical Windows Image Commands}
These commands require the separately retained Windows image. Set the
PowerShell variable to its actual path first.
\begin{lstlisting}
$IMAGE = "path/to/NTFS_Forensic_Baseline_ADS_Slack_Working_RAW.dd"
py code/ntfs_image_reader.py $IMAGE --show-record 41 --show-record 58
py code/ntfs_recovery.py $IMAGE --record 58 --stream CapstoneEvidence --output new_lab_ads.bin
py code/ntfs_recovery.py $IMAGE --record 41 --slack-length 25 --output new_lab_slack.bin
\end{lstlisting}
For a saved verification decision, add these options with the original values:
\begin{lstlisting}
--report NEW_REPORT.json
--expected-length ORIGINAL_BYTE_COUNT
--expected-sha256 ORIGINAL_PAYLOAD_HASH
--expected-image-sha256 ORIGINAL_IMAGE_HASH
\end{lstlisting}
Appendix B records the hashes. Output parents must exist; destinations must
be new. Full-image hashing takes longer than locating content.

\section{Interpreting Decisions}
PASS requires the expected payload length and original hash to match, with
observed source preservation. FAIL indicates a mismatched expectation or
changed source. NOT\_VERIFIED means an original expectation was omitted.
Input or recovery errors are distinct from a failed content comparison.
The reporting CLI returns 0 for PASS, 1 for FAIL or NOT\_VERIFIED, and 2 for
input/recovery errors. Diagnostic or partial files after an error must not
be treated as completed evidence.

\chapter{Complete Raw Test Results}
\section{Marker and Image Hash Register}
\noindent\textbf{Working-image SHA-256:}\\
\workinghash

\medskip\noindent\textbf{ADS marker, 36 bytes:}\\
\file{CAPSTONE_ADS_TEST_20261001_Team3_001}\\
\texttt{E3CB29B996899849E804A8FD7DB33F7CC}\\
\texttt{ECBCB2569005353BD082D320D0D997E}

\medskip\noindent\textbf{Slack marker, 25 bytes:}\\
\file{TEAM3_SLACK_TEST_20261008}\\
\texttt{3BF8F582FF237AD7F739D39EDCF54CB35E}\\
\texttt{0E6FDFD75AD47CA55F3598629390F2}

\medskip\noindent Windows slack location: record 41, absolute byte offset
16,049,277; available tail 2,947 bytes; selected payload 25 bytes.
These markers and their hashes precede the later recovery comparisons.
The portable fixture uses the same marker strings but a different image
and geometry; its image hash must not be confused with the Windows image.

\section{Baseline File Hash Register}
The November 15 recovery recorded the hashes below. On November 25, Jeff confirmed that the original-file hashes match all eight recorded values. Original values are reproduced on that team confirmation; original command output was not attached. The dated confirmation is stored in reports/evidence/final\_report/baseline\_hash\_confirmation.json.
\begin{longtable}{p{1.6in}p{4.4in}}
\caption{Original and recovered SHA-256 values confirmed equal for all eight files.}\\
\toprule File & Original and recovered SHA-256\\\midrule\endfirsthead
\toprule File & Original and recovered SHA-256\\\midrule\endhead
Test.txt & \texttt{D68EC557336622896727E44ED0031366}\newline\texttt{62FCD67CE8F9F57B23E70C0E09ADA797}\\[0.6em]
Test2.txt & \texttt{BA7F25F70009A6D45A9EF88C857C9C35}\newline\texttt{7856AA6D1201486DE9070522CFAEE64E}\\[0.6em]
Testfile.docx & \texttt{84A8A7AC155C36A5F4F9196A412D8DB1}\newline\texttt{6A6F12BD10794FEC362542E562DBE347}\\[0.6em]
SunflowerTest.jpg & \texttt{658C1F66CE1C2D05B51F4E2004498122}\newline\texttt{8AAE39BE33B14737B5ADD448378B9E7C}\\[0.6em]
Thetestfile.pdf & \texttt{61FCC2B322094257A6B845E15D7AEFF9}\newline\texttt{9564E83260EA413DB05180DD8A0E3C8C}\\[0.6em]
WaterfallTest.png & \texttt{19CE26BC665962B8B6B766C48EFFC867}\newline\texttt{6D1191B7B46CACD165418892ED42AA08}\\[0.6em]
HTMLtestdoc.html & \texttt{815246781F0ADE404A0FECD2531CB027}\newline\texttt{AB7F6C1177868967BC7CCA47D88AE0FB}\\[0.6em]
JSON\_test\_File.json & \texttt{65EAF0824C357D4FC9955E75CA29CF3F}\newline\texttt{6476DCB0374241DAB08C3714991E8D34}\\[0.6em]
\bottomrule\end{longtable}
\section{All Synthetic Payload Samples}
Each row represents a separately generated payload. Clean recovery and the direct-read baseline matched exactly in every row. Each changed case contains one deliberately flipped bit and a failed original-payload hash comparison.
\begin{longtable}{rrrrrr}
\caption{Complete 50-sample synthetic campaign results.}\\
\toprule ID & Bytes & Clean BER \% & Changed BER \% & Exact baseline & Hash failure\\\midrule\endfirsthead
\toprule ID & Bytes & Clean BER \% & Changed BER \% & Exact baseline & Hash failure\\\midrule\endhead
1 & 1 & 0 & 12.500000 & Yes & Yes\\
2 & 25 & 0 & 0.500000 & Yes & Yes\\
3 & 100 & 0 & 0.125000 & Yes & Yes\\
4 & 200 & 0 & 0.062500 & Yes & Yes\\
5 & 324 & 0 & 0.038580 & Yes & Yes\\
6 & 1 & 0 & 12.500000 & Yes & Yes\\
7 & 25 & 0 & 0.500000 & Yes & Yes\\
8 & 100 & 0 & 0.125000 & Yes & Yes\\
9 & 200 & 0 & 0.062500 & Yes & Yes\\
10 & 324 & 0 & 0.038580 & Yes & Yes\\
11 & 1 & 0 & 12.500000 & Yes & Yes\\
12 & 25 & 0 & 0.500000 & Yes & Yes\\
13 & 100 & 0 & 0.125000 & Yes & Yes\\
14 & 200 & 0 & 0.062500 & Yes & Yes\\
15 & 324 & 0 & 0.038580 & Yes & Yes\\
16 & 1 & 0 & 12.500000 & Yes & Yes\\
17 & 25 & 0 & 0.500000 & Yes & Yes\\
18 & 100 & 0 & 0.125000 & Yes & Yes\\
19 & 200 & 0 & 0.062500 & Yes & Yes\\
20 & 324 & 0 & 0.038580 & Yes & Yes\\
21 & 1 & 0 & 12.500000 & Yes & Yes\\
22 & 25 & 0 & 0.500000 & Yes & Yes\\
23 & 100 & 0 & 0.125000 & Yes & Yes\\
24 & 200 & 0 & 0.062500 & Yes & Yes\\
25 & 324 & 0 & 0.038580 & Yes & Yes\\
26 & 1 & 0 & 12.500000 & Yes & Yes\\
27 & 25 & 0 & 0.500000 & Yes & Yes\\
28 & 100 & 0 & 0.125000 & Yes & Yes\\
29 & 200 & 0 & 0.062500 & Yes & Yes\\
30 & 324 & 0 & 0.038580 & Yes & Yes\\
31 & 1 & 0 & 12.500000 & Yes & Yes\\
32 & 25 & 0 & 0.500000 & Yes & Yes\\
33 & 100 & 0 & 0.125000 & Yes & Yes\\
34 & 200 & 0 & 0.062500 & Yes & Yes\\
35 & 324 & 0 & 0.038580 & Yes & Yes\\
36 & 1 & 0 & 12.500000 & Yes & Yes\\
37 & 25 & 0 & 0.500000 & Yes & Yes\\
38 & 100 & 0 & 0.125000 & Yes & Yes\\
39 & 200 & 0 & 0.062500 & Yes & Yes\\
40 & 324 & 0 & 0.038580 & Yes & Yes\\
41 & 1 & 0 & 12.500000 & Yes & Yes\\
42 & 25 & 0 & 0.500000 & Yes & Yes\\
43 & 100 & 0 & 0.125000 & Yes & Yes\\
44 & 200 & 0 & 0.062500 & Yes & Yes\\
45 & 324 & 0 & 0.038580 & Yes & Yes\\
46 & 1 & 0 & 12.500000 & Yes & Yes\\
47 & 25 & 0 & 0.500000 & Yes & Yes\\
48 & 100 & 0 & 0.125000 & Yes & Yes\\
49 & 200 & 0 & 0.062500 & Yes & Yes\\
50 & 324 & 0 & 0.038580 & Yes & Yes\\
\bottomrule\end{longtable}
\section{Release Regression Output}
The following complete output is the November 24 clean-copy release test run. The seven acceptance checks are recorded in the accompanying verification JSON.
\lstinputlisting[basicstyle=\ttfamily\scriptsize]{../evidence/final_code/regression.txt}
\section{Acceptance Summary}
\lstinputlisting[basicstyle=\ttfamily\scriptsize]{../evidence/final_code/verification.json}


\section{Detection Campaign: All 100 Decisions}
The following raw CSV retains every construction label, prediction, outcome,
and source-preservation decision from November 25. The header identifies the
columns. True means a present hidden payload, a flag, or unchanged source as
appropriate. Interpret the two screening groups separately using Chapter 5's
denominators; this file is not a recovery-hash comparison.
\lstinputlisting[basicstyle=\ttfamily\scriptsize]{../evidence/final_detection/results.csv}

\section{Raw Artifact Locations}
The repository retains the full historical image and slack validation JSON
under \file{reports/evidence/week13/}, release output under
\file{reports/evidence/final_code/}, and all 50 campaign rows with timing
columns under \file{reports/evidence/final_report/results.csv}.
The campaign generator recreates individual originals, images, outputs,
and reports; those generated binaries need not be stored in Git. The
historical evidence catalog contains additional image-state hashes,
screenshots, and hexdumps.

\chapter{Development Log and Contribution Record}
\section{Version History Summary}
The table below was generated from non-merge commits reachable from the
report branch's starting point, through November 24. Merge commits are
excluded to avoid counting integration as another implementation change.
Jeff's historical identities are grouped together; the Lucas identity is
\file{Lcurtis1Broncos}. Counts are evidence of recorded activity, not
measurements of effort, quality, or a prescribed individual grade.
The underlying dated commit list is retained in
\file{reports/evidence/final_report/commit_history.txt}.

\begin{longtable}{lrrr}
\caption{Non-merge commits by week and recorded author identity.}\\
\toprule Week beginning & Jeff identities & Lucas identity & Other\\\midrule\endfirsthead
\toprule Week beginning & Jeff identities & Lucas identity & Other\\\midrule\endhead
2026-08-17 & 3 & 6 & 0\\
2026-08-24 & 1 & 2 & 0\\
2026-08-31 & 13 & 19 & 0\\
2026-09-07 & 5 & 7 & 0\\
2026-09-14 & 4 & 4 & 0\\
2026-09-21 & 2 & 1 & 0\\
2026-09-28 & 1 & 6 & 0\\
2026-10-05 & 1 & 0 & 0\\
2026-10-12 & 1 & 0 & 0\\
2026-10-19 & 1 & 1 & 0\\
2026-10-26 & 2 & 5 & 0\\
2026-11-02 & 1 & 5 & 0\\
2026-11-09 & 2 & 16 & 0\\
2026-11-16 & 0 & 4 & 0\\
2026-11-23 & 1 & 0 & 0\\
\bottomrule\end{longtable}


\section{Contribution Summary}
Lucas prepared the dedicated VirtualBox laboratory, supplied image exports
and known-file information, documented the manual ADS procedure, and added
ADS/slack workflow records, an evidence catalog, and the initial final-report
and evidence-log drafts. His repository activity also records review,
documentation updates, and integration of weekly work.

Jeff implemented the boot and MFT parser, attribute and run interpretation,
controlled slack-copy tool, recovery support and verification reporting,
automated tests, and release packaging. He recorded parser and payload
validation evidence and assembled weekly documentation. Jeff also prepared
the November 24 report campaign and expanded report.

The meeting records and journals provide context alongside commits. The
semester plan's midterm interval was adjusted after the actual deadline
required submission on October 11. The following week was used for validation
and planning. Later work progressed through recovery, ADS checks, slack
checks, consolidated testing, and final evidence/report preparation. This
sequence is visible in the dated source history and weekly evidence.

\section{Repository Contents}
The tagged code package includes source modules, tests and dataset generators,
README instructions, design notes, evidence, and weekly documentation. The
report branch adds the final report source, figures, campaign generator, and
raw campaign tables. Large original images and personal lesson exports
remain outside Git. For the code assignment the immutable submission
reference is \file{v1.0-final}; the signed final report is a separate artifact.
